%=====================================================================
\chapter{Cardinality, Participation and Weak Entities}\label{ch:cardinality}
%=====================================================================
\locator{2}{Part II --- Designing It}

\begin{outcomes}
By the end of this chapter you will be able to:
\begin{tight}
  \item \textbf{Determine} the cardinality ratio of a relationship --- 1:1,
        1:N or M:N --- by asking two questions in a fixed order.
  \item \textbf{Distinguish} total from partial participation, and
        \textbf{state} what each one forbids.
  \item \textbf{Write} a relationship's constraints in $(\min,\max)$ notation
        and \textbf{convert} between that and the ratio-plus-participation
        form.
  \item \textbf{Recognise} a weak entity type, and \textbf{identify} its owner,
        its identifying relationship and its partial key.
  \item \textbf{Read} the same diagram in Chen and crow's-foot notation.
\end{tight}
\end{outcomes}

\begin{prereq}
\chref{ermodel}: entities, attributes, keys, relationships and degree. You
will extend the diagram you drew in that chapter's laboratory, so have it with
you.
\end{prereq}

\begin{keyterms}
cardinality ratio $\bullet$ one-to-one $\bullet$ one-to-many $\bullet$
many-to-many $\bullet$ participation constraint $\bullet$ total participation
$\bullet$ partial participation $\bullet$ existence dependency $\bullet$
structural constraint $\bullet$ $(\min,\max)$ notation $\bullet$ weak entity
type $\bullet$ strong entity type $\bullet$ owner entity $\bullet$ identifying
relationship $\bullet$ partial key $\bullet$ discriminator $\bullet$
crow's-foot notation
\end{keyterms}

\section{A Diagram Without Cardinalities Says Almost Nothing}

The diagram from \chref{ermodel} says students enrol in sections. It does not
say whether a student may take two sections of the same course, whether a
section can run with nobody in it, or whether an instructor must be teaching
something. Those are the facts an examinations officer actually argues about,
and all three are cardinality or participation.

The two constraints answer different questions and are routinely confused, so
they are worth separating before either is defined.

\begin{definitionbox}[Two different questions]
The \term{cardinality ratio} answers \emph{how many}: given one entity on this
side, how many on the other side may it be related to --- one, or more than
one?

\smallskip
The \term{participation constraint} answers \emph{must it at all}: can an
entity on this side exist without taking part in the relationship?

\smallskip
They are independent. Every combination occurs, and a relationship is not
specified until both have been answered for \emph{both} ends --- four answers
in all.
\end{definitionbox}

\section{Cardinality, by Asking Two Questions}

\begin{definitionbox}[The three ratios]
\term{One-to-one (1:1)}: each entity on either side relates to at most one on
the other. A department has one head; a person heads at most one department.

\smallskip
\term{One-to-many (1:N)}: one entity on the ``one'' side relates to many on
the ``many'' side, but each of those relates back to at most one. A department
offers many courses; each course belongs to one department.

\smallskip
\term{Many-to-many (M:N)}: entities on both sides may relate to many. A student
enrols in many sections; a section holds many students.
\end{definitionbox}

\begin{conceptbox}[The procedure, and it is only two questions]
Never guess the ratio. Ask these, in this order, out loud, with a concrete
entity in each:

\begin{tight}
  \item \textbf{``One \emph{student} --- how many \emph{sections}?''}
        \quad One, or many?
  \item \textbf{``One \emph{section} --- how many \emph{students}?''}
        \quad One, or many?
\end{tight}

\smallskip
Two answers, four combinations, three ratios:

\medskip
\begin{tabular}{@{}L{3.3cm}L{3.3cm}L{3.0cm}L{4.1cm}@{}}
\textbf{Q1} & \textbf{Q2} & \textbf{Ratio} & \textbf{Example} \\
one  & one  & 1:1 & department --- head \\
many & one  & 1:N & department --- course \\
one  & many & 1:N & \emph{the same, read backwards} \\
many & many & M:N & student --- section \\
\end{tabular}

\smallskip
\textbf{Rows 2 and 3 are the same relationship.} ``1:N'' is meaningless until
you say which end is the one. Always write it as \ent{department}
1---N \ent{course}, never as a bare ``1:N''.
\end{conceptbox}

\begin{examplebox}[The question that catches everybody]
\emph{``A student enrols in many sections. A section has many students.
Therefore M:N.''} --- correct. Now:

\medskip
\emph{``A student enrols in many \textbf{courses}. A course has many students.
Therefore M:N.''} --- also correct, and \textbf{it is the wrong relationship}.

\smallskip
A student does not enrol in a course. They enrol in a \emph{section} of a
course, in a particular semester, taught by a particular instructor, in a
particular room. Model the relationship to \ent{course} and you can never
record which section somebody attended, which instructor taught them, or that
they failed in Fall and passed in Spring.

\smallskip
\textbf{The lesson is not about cardinality.} It is that a cardinality
question asked about the wrong pair of entity types gives a confident,
well-formed, useless answer. Settle what the entities are
(\chref{ermodel}) before asking how many.
\end{examplebox}

\section{Participation: Must It Take Part at All?}

\begin{definitionbox}[Total and partial participation]
\term{Total participation} of an entity type in a relationship means
\emph{every} entity of that type must take part. It is drawn as a double line
and is also called an \term{existence dependency}: a section cannot exist
without a course.

\smallskip
\term{Partial participation} means some entities may take part and others need
not. A single line. A student may be enrolled in nothing --- newly admitted,
on leave --- and is still a student.

\smallskip
Total participation is the stronger claim, and therefore the one to be
sceptical about: it forbids something, and somebody will eventually need the
thing it forbids.
\end{definitionbox}

\begin{alertbox}[Total participation is a promise about Tuesday afternoon]
``Every section must be taught by an instructor'' sounds obviously true, and it
is how timetables are published. Then:

\begin{tight}
  \item A section is created in June and the instructor is appointed in August.
        For two months it has none.
  \item An instructor resigns mid-semester. For a week the section has none.
  \item A section is cancelled but its enrolment records must be kept.
\end{tight}

\smallskip
If participation is total, the database \emph{forbids} all three, and the clerk
who meets them invents a fake instructor called ``TBA'' --- which now appears
in the staff list, the workload report and the payroll query. A constraint
that gets worked around is worse than one that was never there, because the
data now lies and nothing says so.

\smallskip
\textbf{The decision rule.} Make participation total when the entity is
genuinely meaningless without the relationship --- a section without a course
is not a thing at all. Make it partial when the relationship is merely
\emph{usual}. ``Always'' and ``usually'' are different words and the database
can tell.
\end{alertbox}

\begin{definitionbox}[$(\min,\max)$ notation]
Chen's two constraints can be written as one pair on each side. Against each
entity type's line, $(\min,\max)$ gives the smallest and largest number of
relationship instances \emph{one entity of that type} may take part in.

\smallskip
$\min = 0$ is partial participation; $\min \ge 1$ is total. $\max = 1$ is the
``one'' side; $\max = N$ is the ``many'' side.

\smallskip
\textbf{The notation is read against the entity it is nearest}, which is the
opposite of where the 1 and N labels go in classic Chen notation. This single
difference is responsible for most misread diagrams, and is why many
examiners insist on one or the other but not both.
\end{definitionbox}

\begin{worked}[The university's four relationships, specified completely]
Four answers each: ratio from both sides, participation from both sides.

\medskip
\textbf{1. \ent{department} offers \ent{course}}
\begin{tight}
  \item One department --- many courses. One course --- one department.
        \textbf{1:N.}
  \item Must every department offer a course? A new department may not yet.
        \textbf{Partial} (0,N).
  \item Must every course belong to a department? Yes; it is how it is
        numbered. \textbf{Total} (1,1).
\end{tight}
$\Rightarrow$ \ent{department}~(0,N)~---~\ent{offers}~---~(1,1)~\ent{course}

\medskip
\textbf{2. \ent{instructor} teaches \ent{section}}
\begin{tight}
  \item One instructor --- many sections. One section --- one instructor
        (this university's rule; some allow co-teaching, and then it is M:N).
        \textbf{1:N.}
  \item Must every instructor teach? On sabbatical, no. \textbf{Partial} (0,N).
  \item Must every section have an instructor? See the alert box above:
        \textbf{partial} (0,1), deliberately.
\end{tight}

\medskip
\textbf{3. \ent{student} enrols in \ent{section}}
\begin{tight}
  \item Many each way. \textbf{M:N}, and it carries \texttt{grade}.
  \item Must every student be enrolled in something? No. \textbf{Partial} (0,N).
  \item Must every section have students? A section with nobody in it is
        cancelled, not impossible. \textbf{Partial} (0,N).
\end{tight}

\medskip
\textbf{4. \ent{course} requires \ent{course}} (recursive)
\begin{tight}
  \item A course may require several; a course may be required by several.
        \textbf{M:N}, both ends partial (0,N) --- most courses have no
        prerequisites and are a prerequisite for nothing.
\end{tight}

\medskip
\textbf{What the four tell you.} Three of the four are partial on both ends.
Beginners' diagrams are covered in double lines, because every relationship
\emph{sounds} compulsory when you describe the normal case. The normal case is
not the constraint.
\end{worked}

\section{Weak Entities}

\begin{definitionbox}[Weak entity type, owner, identifying relationship, partial key]
A \term{weak entity type} has no key of its own. Its entities are
distinguishable only in combination with the entity that owns them.

\smallskip
The \term{owner} (or identifying) entity type is the strong one it depends on.
The \term{identifying relationship} joins them, drawn as a double diamond, and
the weak entity's participation in it is always \emph{total}.

\smallskip
The \term{partial key} (or \term{discriminator}) is the attribute that tells
one weak entity from another \emph{belonging to the same owner}. It is
underlined with a dashed line.

\smallskip
The full key of a weak entity is \textbf{the owner's primary key plus the
partial key}.
\end{definitionbox}

\begin{examplebox}[\ent{section} is weak, and this is why]
Sections are numbered 1, 2, 3 within a course. Section 1 of CS-103 and section
1 of CS-201 are different sections with the same number.

\begin{tight}
  \item \texttt{section\_no} alone is not a key --- dozens of sections are
        numbered 1.
  \item \texttt{section\_no} \emph{within one course} is unique. That makes it
        a \textbf{partial key}.
  \item A section cannot exist without a course. \textbf{Total participation}
        in the identifying relationship.
  \item Its real key is $\{$\texttt{course\_code}, \texttt{semester},
        \texttt{section\_no}$\}$ --- the owner's key, plus the semester,
        plus the discriminator.
\end{tight}

\smallskip
\relation{section}{\relkey{course\_code}, \relkey{semester},
\relkey{section\_no}, room, slot, emp\_id}

\smallskip
\textbf{The alternative, and why it is not better.} You could give
\ent{section} a surrogate \texttt{section\_id} and make it strong. Many real
systems do, and \chref{mapping} shows the mechanics. It does not make the
entity any less dependent --- it only hides the dependency from the key, and
the uniqueness of (course, semester, number) must still be declared or the
same section can be created twice.
\end{examplebox}

\dsfig{%
\begin{tikzpicture}[
  ent/.style={draw=primary, line width=1.1pt, fill=primary!8,
              text=primary!60!black, font=\scriptsize\bfseries,
              minimum width=1.9cm, minimum height=0.68cm},
  weak/.style={ent, draw=purplehl, fill=purplehl!8, text=purplehl!65!black,
               double, double distance=1.2pt},
  rel/.style={draw=greenhl, line width=0.9pt, fill=greenhl!8,
              text=greenhl!55!black, font=\scriptsize\bfseries,
              diamond, aspect=2.2, inner sep=0.5pt},
  irel/.style={rel, draw=purplehl, fill=purplehl!8, text=purplehl!65!black,
               double, double distance=1.2pt},
  one/.style={draw=neutral!60, line width=0.7pt},
  two/.style={draw=neutral!60, line width=0.7pt, double,
              double distance=1.1pt},
  mm/.style={font=\scriptsize\bfseries, text=accent, inner sep=1.2pt},
  lbl/.style={font=\scriptsize, text=neutral}
]
  % Centres 3.2 apart on both rows. The first draft used 2.3 and the TEACHES
  % diamond printed on top of the INSTRUCTOR box -- a diamond of aspect 2.2 is
  % wider than it looks, and three of the (min,max) labels landed on boxes for
  % the same reason. The labels now sit at pos=0.22, which also puts each one
  % nearest the entity it is read against, as the definition requires.
  \node[ent]  (dep) at (-7.6, 1.7) {DEPARTMENT};
  \node[rel]  (off) at (-3.8, 1.7) {OFFERS};
  \node[ent]  (crs) at ( 0.0, 1.7) {COURSE};
  \node[rel]  (req) at ( 0.0, 3.9) {REQUIRES};
  \node[irel] (has) at ( 0.0, 0.0) {HAS};
  \node[weak] (sec) at ( 0.0,-1.9) {SECTION};
  \node[ent]  (stu) at (-7.6,-1.9) {STUDENT};
  \node[rel]  (enr) at (-3.8,-1.9) {ENROLS};
  \node[rel]  (tch) at ( 3.8,-1.9) {TEACHES};
  \node[ent]  (ins) at ( 7.6,-1.9) {INSTRUCTOR};

  \draw[one] (dep) -- node[mm, above=1pt, pos=0.35] {(0,N)} (off);
  \draw[two] (off) -- node[mm, above=1pt, pos=0.65] {(1,1)} (crs);
  \draw[one] (crs) -- node[mm, right=1pt, pos=0.3] {(0,N)} (has);
  \draw[two] (has) -- node[mm, right=1pt, pos=0.7] {(1,1)} (sec);
  \draw[one] (stu) -- node[mm, above=1pt, pos=0.35] {(0,N)} (enr);
  \draw[one] (enr) -- node[mm, above=1pt, pos=0.65] {(0,N)} (sec);
  \draw[one] (sec) -- node[mm, above=1pt, pos=0.35] {(0,1)} (tch);
  \draw[one] (tch) -- node[mm, above=1pt, pos=0.65] {(0,N)} (ins);

  % The recursive relationship, as a loop above COURSE so that both of its
  % ends are visible and can carry role names. Drawn flat on the right of the
  % course box it read as a line passing over the top to nothing.
  \draw[one] (crs.north west) .. controls (-1.9,2.6) and (-1.9,3.9)
             .. (req.west);
  \draw[one] (crs.north east) .. controls ( 1.9,2.6) and ( 1.9,3.9)
             .. (req.east);
  \node[lbl, anchor=east, font=\scriptsize\itshape] at (-1.25,3.15)
       {is prerequisite for};
  \node[lbl, anchor=west, font=\scriptsize\itshape] at ( 1.25,3.15)
       {depends on};
  \node[mm, anchor=south east] at (-0.35,2.15) {(0,N)};
  \node[mm, anchor=south west] at ( 0.35,2.15) {(0,N)};

  \node[lbl, text=purplehl!70!black, text width=4.2cm, align=left,
        anchor=west] at (0.95,0.0)
       {double diamond and double line: identifying, and always total};
  \node[lbl, text=purplehl!70!black, anchor=north, text width=3.4cm,
        align=center] at (0.0,-2.45)
       {weak: no key of its own};
  % No explicit \\ here. An align=right node whose forced line does not fill
  % its text width is reported as an underfull hbox, twice, once per line.
  % Letting it wrap on its own costs nothing and the gate stays clean.
  \node[lbl, text=accent, anchor=east, align=right, text width=3.8cm]
       at (-4.6,0.3) {$(\min,\max)$ is read against the entity it touches};
\end{tikzpicture}}%
{The university database, fully constrained. Only two lines in the whole
diagram are double --- a course must belong to a department and a section must
belong to a course. Everything else is partial, which is the usual
proportion.}{ercomplete}

\begin{conceptbox}[Chen or crow's foot? Both, and here is the translation]
Textbooks and examination papers use Chen. Every tool you will meet at work
uses crow's foot, because it fits on a screen. They say the same things.

\medskip
\begin{tabular}{@{}L{4.4cm}L{4.6cm}L{4.7cm}@{}}
\textbf{Meaning} & \textbf{Chen} & \textbf{Crow's foot} \\
relationship & a diamond between boxes & a line between boxes \\
``many'' end & N on the line & the three-pronged foot \\
``one'' end & 1 on the line & a single cross-stroke \\
total participation & double line & a cross-stroke nearest the box \\
partial participation & single line & a circle nearest the box \\
attributes & ellipses around the box & listed inside the box \\
weak entity & double rectangle & usually not distinguished \\
M:N relationship & one diamond & \textbf{cannot be drawn} --- it is
resolved into two 1:N at once \\
\end{tabular}

\smallskip
\textbf{The last row is the real difference.} Crow's-foot notation has no way
to draw a many-to-many relationship, so a tool forces you to resolve it into a
third box before you have decided to. Chen lets you keep the design honest
until \chref{mapping}, where the resolution is a rule with a reason rather
than something the drawing program did to you.
\end{conceptbox}

\begin{pitfall}
\begin{tight}
  \item \textbf{Writing ``1:N'' without saying which end is the one.} The
        label alone is ambiguous; name the entity types.
  \item \textbf{Confusing cardinality with participation.} ``Every section has
        one instructor'' states both --- one is the ratio, every is the
        participation --- and people hear only the first.
  \item \textbf{Double lines everywhere.} Every relationship sounds compulsory
        when the normal case is described. Three of this database's four are
        partial on both ends.
  \item \textbf{Reading $(\min,\max)$ against the far entity.} It is read
        against the entity it touches, which is the opposite of the 1/N labels.
  \item \textbf{Asking the cardinality of the wrong pair.} Student-to-course is
        M:N and is the wrong relationship; student-to-section is the right one.
  \item \textbf{Giving a weak entity a surrogate key and calling it strong.}
        The dependency does not go away; it just stops being visible, and the
        uniqueness of the real key must still be declared.
\end{tight}
\end{pitfall}

\begin{lab}[Constrain your diagram, then try to break it]
Take the diagram from \chref{ermodel}'s laboratory. Still no database.

\begin{lstlisting}
# --- 1. Four answers per relationship -------------------------------
    For EVERY relationship, write all four:
        one <A> -- how many <B>?        one <B> -- how many <A>?
        must every <A> take part?       must every <B> take part?
    Four answers. Not three.

# --- 2. Mark the ratio ----------------------------------------------
    1:1, 1:N or M:N, naming the entity types -- not a bare "1:N".

# --- 3. Mark participation ------------------------------------------
    Double line for total, single for partial.

# --- 4. Rewrite it in (min,max) -------------------------------------
    Every line gets a pair at each end. Check each against step 1;
    they must agree, and if they do not, step 1 was guessed.

# --- 5. Find the weak entity ----------------------------------------
    There is at least one. Name its owner, its identifying
    relationship, its partial key and its full key.

# --- 6. Attack every double line ------------------------------------
    For EACH total participation, invent a real Tuesday on which the
    university needs to break it. If you can, make it partial.
    Write the Tuesday down next to the line; it is the evidence.

# --- 7. Attack every "one" ------------------------------------------
    For each 1 end, ask what happens if the university someday allows
    two. Co-teaching. Joint degrees. Two heads of department.
    Which of these would cost a schema change, and which would not?

# --- 8. Read it back ------------------------------------------------
    Turn the finished diagram into ten English sentences, one per
    constraint. Give them to the person who has the paragraph from
    Chapter 5's lab. Every disagreement is a decision you had not
    actually made.
\end{lstlisting}

\smallskip
\textbf{What to notice.} Steps 6 and 7 usually remove more constraints than
they confirm, and that is the expected result rather than a sign you did it
wrong. A design is not improved by forbidding more; it is improved by
forbidding exactly what the university actually forbids.
\end{lab}

\begin{checkpoint}
\begin{tightnum}
  \item A relationship is drawn with a double line at the \ent{section} end
        of \ent{teaches}. State in one sentence what the database will
        refuse to do, and give a realistic situation in which somebody needs it
        to do exactly that.
  \item \ent{enrolment} has $(0,N)$ at the \ent{student} end and $(0,N)$
        at the \ent{section} end. Rewrite this as a cardinality ratio plus
        two participation constraints, and say what would change if the
        \ent{section} end became $(1,N)$.
  \item Explain why \ent{section} is a weak entity type, name its partial
        key and its full key, and then say what you would have to declare
        instead if you chose to give it a surrogate \texttt{section\_id} ---
        and what breaks if you forget.
\end{tightnum}
\end{checkpoint}

\begin{chaptersummary}
\begin{tight}
  \item A relationship needs \textbf{four} answers: the ratio from each side
        and the participation from each side. Fewer than four is an
        unspecified relationship.
  \item The \textbf{cardinality ratio} is found by two questions --- ``one A,
        how many B?'' and ``one B, how many A?'' --- giving 1:1, 1:N or M:N.
        Always name which end is the one.
  \item \textbf{Total participation} (double line) means every entity must take
        part; \textbf{partial} means it need not. Total is the stronger claim
        and deserves the scepticism.
  \item A constraint that gets worked around is worse than no constraint: the
        data then contains a fake row and nothing marks it as fake.
  \item $(\min,\max)$ carries both facts at once: $\min=0$ is partial,
        $\min\ge 1$ is total, $\max=1$ is the one side. \textbf{It is read
        against the entity it touches}, unlike the 1/N labels.
  \item A \textbf{weak entity type} has no key of its own. It has an
        \textbf{owner}, an \textbf{identifying relationship} (double diamond,
        always total) and a \textbf{partial key}; its full key is the owner's
        key plus the partial key.
  \item \ent{section} is weak: section 1 exists in every course, and only
        (course, semester, number) identifies one.
  \item Crow's-foot notation says the same things more compactly but
        \textbf{cannot draw an M:N relationship}, so a tool resolves it for you
        before you have decided to.
  \item Asking a cardinality question about the wrong pair of entity types
        gives a confident and useless answer. Student-to-\emph{course} is the
        classic case.
\end{tight}
\end{chaptersummary}

\begin{reviewq}
  \item \textbf{(MCQ)} A double line between an entity type and a relationship
        means: (a)~many-to-many; (b)~total participation; (c)~a weak entity;
        (d)~an identifying relationship.
  \item \textbf{(MCQ)} In $(\min,\max)$ notation, $(1,1)$ against
        \ent{course} on the \ent{offers} line means: (a)~a department
        offers exactly one course; (b)~a course belongs to exactly one
        department; (c)~there is one department; (d)~the relationship is 1:1.
  \item \textbf{(Short)} Give an example of a 1:1 relationship in a university
        and say which end you would make total, with a reason.
  \item \textbf{(Short)} Why is the participation of a weak entity type in its
        identifying relationship always total? Answer from the definition of a
        weak entity.
  \item \textbf{(Short)} A tool that draws crow's-foot diagrams will not let
        you record a many-to-many relationship. What does it make you do
        instead, and what is lost by doing it at that moment rather than in
        \chref{mapping}?
  \item \textbf{(Applied)} A hospital records patients, doctors, appointments
        and prescriptions. Give all four answers for every relationship you
        identify, state which entity types are weak and why, and defend each
        total participation with a specific situation that would break it if
        you were wrong.
  \item \textbf{(Applied)} Your university decides to allow co-teaching: two
        instructors may share a section. State exactly which constraint in
        Figure~\ref{fig:ercomplete} changes, what the relationship becomes,
        what new attributes it may now need, and --- looking ahead to
        \chref{mapping} --- why this change is cheap on paper and expensive
        once the database is full.
\end{reviewq}
